\subsection{Sémantická analýza}
\label{sec:analyzer}

Sémantická analýza (\texttt{src/analyzer.rs}) ověřuje, že program používá pouze
existující slova a~že uživatelské definice neporušují pravidla jazyka.
V~kompilátoru \\Roth{} je analýza založena na udržování tří tabulek:
\texttt{builtin\_words}, \texttt{defined\_words} a~\texttt{defined\_variables}.

Při inicializaci analyzátoru je tabulka vestavěných slov naplněna seznamem
primitivních operací (aritmetika, práce se zásobníkem, I/O, řízení toku,
práce s~pamětí). Vestavěná slova jsou považována za součást definice jazyka a~nelze
je redefinovat.

Analýza prochází AST rekurzivně:
\begin{itemize}
	\item při uzlu \texttt{Definition} ověří, že jméno není vestavěné, následně
	      přidá slovo mezi definovaná ještě \emph{před} analýzou těla;
	      tím je umožněna rekurze pomocí \texttt{RECURSE},
	\item při uzlu \texttt{Word} ověří, že se jedná o~vestavěné slovo, uživatelsky
	      definované slovo, nebo proměnnou,
	\item při \texttt{VariableDeclaration} registruje proměnnou v tabulce.
\end{itemize}

V~režimu REPL je tato fáze rozšířena o~kontext předchozích vstupů: před voláním
\texttt{analyze} jsou do analyzátoru doplněny již definované slova a~proměnné,
aby bylo možné na ně navazovat (viz sekce~\ref{sec:repl}).
